\section[The addition formula for the $\tau$-function of the KP hierarchy]{The addition formula for the $\boldsymbol{\tau}$-function of the KP hierarchy} \label{section2}

Let
\begin{gather*}
[\alpha]=\left(\alpha,\frac{\alpha^2}{2},\frac{\alpha^3}{3},\dots\right),
\qquad
\xi(t,\lambda)=\sum_{n=1}^{\infty}t_n\lambda^n,
\qquad
t=(t_1,t_2,t_3,\dots).
\end{gather*}
The KP hierarchy is a~system of equations for a~function $\tau(t)$ \cite{DJKM1,JM1, MJD1} given by
\begin{gather}
\oint e^{\xi(t'-t,\lambda)}\tau\big(t'-\big[\lambda^{-1}\big]\big)\tau\big(t+\big[\lambda^{-1}\big]\big)\frac{d\lambda}{2\pi i}=0.
\label{eq:1}
\end{gather}
Here $\oint$ means a~formal algebraic operator extracting the coef\/f\/icient of $z^{-1}$ of Laurent series:
\begin{gather*}%\label{eq:2}
\oint\frac{dz}{2\pi i}\sum_{n=-\infty}^{\infty}a_n z^n=a_{-1}.
\end{gather*}
Set $t=x+y$, $t'=x-y$.
Then~\eqref{eq:1} becomes
\begin{gather}
\oint e^{-2\xi(y,\lambda)}\tau\big(x-y-\big[\lambda^{-1}\big]\big)\tau\big(x+y+\big[\lambda^{-1}\big]\big)\frac{d\lambda}{2\pi i}=0.
\label{eq:3}
\end{gather}
For an integer $m\ge2$, set
\begin{gather}
y=\frac{1}{2}\left(\sum_{i=1}^{m-1}[\beta_i]-\sum_{i=1}^{m+1}[\alpha_i]\right)
\label{eq:22}
\end{gather}
in~\eqref{eq:3}. Then it becomes
\begin{gather}
\oint\exp\left({-\xi\left(\displaystyle\sum_{i=1}^{m-1}[\beta_i]-\sum_{i=1}^{m+1}[\alpha_i],\lambda\right)}\right)
\tau\left(x-\frac{1}{2}\left(\sum_{i=1}^{m-1}[\beta_i]-\sum_{i=1}^{m+1}[\alpha_i]\right)-\big[\lambda^{-1}
\big]\right)\nonumber
\\
\qquad
{}\times\tau\left(x+\frac{1}{2}\left(\sum_{i=1}^{m-1}[\beta_i]-\sum_{i=1}^{m+1}[\alpha_i]\right)+\big[\lambda^{-1}
\big]\right)\frac{d\lambda}{2\pi i}=0.
\label{eq:4}
\end{gather}
By virtue of the identity
\begin{gather*}
\sum_{n=1}^{\infty}\frac{x^n}{n}=-\log(1-x),
\end{gather*}
the exponential factor in~\eqref{eq:4} reduces to a~rational function of $\lambda$, $\alpha_i$, $\beta_i$ as
\begin{gather*}
\exp\left(-\xi\left(\sum_{i=1}^{m-1}[\beta_i]-\sum_{i=1}^{m+1}[\alpha_i],\lambda\right)\right)
=\frac{\prod\limits_{i=1}^{m-1}(1-\beta_i\lambda)}{\prod\limits_{i=1}^{m+1}
(1-\alpha_i\lambda)}.
\end{gather*}
Computing the integral by taking residues at $\lambda=\alpha_i^{-1}$, $1\le i\le m+1$ \cite{M1} and
shifting the variable $x$ as
\begin{gather*}
x\rightarrow x+\frac{1}{2}\left(\sum_{i=1}^{m-1}[\beta_i]+\sum_{i=1}^{m+1}[\alpha_i]\right),
\end{gather*}
we get the following addition formulae of the $\tau$-function~\cite{SS1}
\begin{gather}
\sum_{i=1}^{m+1}(-1)^{i-1}\zeta(x;\beta_1,\dots,\beta_{m-1},\alpha_i)\zeta(x;\alpha_1,\dots,\hat{\alpha}
_i,\dots,\alpha_{m+1})=0,\qquad  m\ge2,
\label{eq:5}
\end{gather}
where
\begin{gather*}
\zeta(x;\alpha_1,\dots,\alpha_{n})=\Delta(\alpha_1,\dots,\alpha_{n})\tau(x+[\alpha_1]+\dots+[\alpha_{n}]),
\\
\Delta(\alpha_1,\dots,\alpha_n)=\prod_{i<j}(\alpha_i-\alpha_j),
\end{gather*}
and $\hat{\alpha}_i$ means that $\alpha_i$ should be removed.
\begin{example}\label{example1}
In the case of $m=2$, we have
\begin{gather}
\alpha_{12}\alpha_{34}\tau(x+[\alpha_1]+[\alpha_2])\tau(x+[\alpha_3]+[\alpha_4])
-\alpha_{13}\alpha_{24}\tau(x+[\alpha_1]+[\alpha_3])\tau(x+[\alpha_2]+[\alpha_4])
\nonumber
\\
\qquad
{}+\alpha_{14}\alpha_{23}\tau(x+[\alpha_1]+[\alpha_4])\tau(x+[\alpha_2]+[\alpha_3])=0,
\label{eq:6}
\end{gather}
where $\alpha_{ij}=\alpha_i-\alpha_j$.
\end{example}

We call~\eqref{eq:6} `the three term equation'.
We have derived~\eqref{eq:6} from~\eqref{eq:1}.
The fact that the converse is true is proved by Takasaki and Takebe~\cite{TT1}.
\begin{theorem}[\cite{TT1}] \label{theorem1} The three term equation~\eqref{eq:6} is equivalent to the KP hierarchy~\eqref{eq:1}.
\end{theorem}

In~\cite{TT1} the theorem is proved by constructing the wave function of the KP-hierarchy.
To do it the dif\/ferential Fay identity, which is a~certain limit of~\eqref{eq:6}, is used.
Here we give an alternative and direct proof of the theorem.
\begin{proposition}\label{proposition1}
The KP hierarchy~\eqref{eq:1} is equivalent to~\eqref{eq:5}.
\end{proposition}
\begin{proof}
It is suf\/f\/icient to prove that if~\eqref{eq:4} holds for any $m\ge2$ and arbitrary
$\alpha_i\neq0$, $\beta_i$, then~\eqref{eq:3} holds.
Set the left hand side of~\eqref{eq:3} to be $F(y)$.
Expand $F(y)$ in $y$ as
\begin{gather*}%\label{eq:7}
F(y)=\sum_\gamma F_\gamma y^\gamma,\qquad  y^\gamma=y_1^{\gamma_1}y_2^{\gamma_2}
\cdots,\qquad \gamma=(\gamma_1,\gamma_2,\dots).
\end{gather*}
We consider the case $\beta_i=0$, $1\le i\le m-1$ in~\eqref{eq:22} and set
\begin{gather*}%\label{eq:8}
y'=\sum_{i=1}^{m+1}[\alpha_i].
\end{gather*}
We prove $F_\gamma=0$ for any $\gamma$ if $F\big({-}\frac{y'}{2}\big)=0$ for any $m\ge2$.
We consider $m$ f\/ixed.
Let us def\/ine the weight of $y_i$ to be $i$ and ${\rm wt}\,y^\gamma=\sum\limits_{i=1}^{\infty}i\gamma_i$.
Decompose $F$ according to weights as
\begin{gather*}
F=F^{(0)}+F^{(1)}+F^{(2)}+\cdots,
%\label{eq:9}
\qquad
F^{(i)}=\sum_{{\rm wt}\,y^\gamma=i}F_\gamma y^\gamma.
%\label{eq:10}
\end{gather*}
We substitute $-\frac{y'}{2}$ to $y$ and get the homogeneous polynomial of degree $i$ of
$\alpha_1,\dots,\alpha_{m+1}$:
\begin{gather*}%\label{eq:11}
F^{(i)}\left(-\frac{y'}{2}\right)=\sum_{\gamma_1+\dots+\gamma_{m+1}=i}b^{(i)}_{\gamma_1\dots\gamma_{m+1}}
\alpha_1^{\gamma_1}\cdots\alpha_{m+1}^{\gamma_{m+1}}.
\end{gather*}
Then $F\big({-}\frac{y'}{2}\big)=0$ is equivalent to $F^{(i)}\big({-}\frac{y'}{2}\big)=0$ for any $i$.
Notice that $i y'_i=\alpha_1^i+\cdots+\alpha_{m+1}^i$ is a~power sum symmetric function.
Therefore $y_1',\dots,y_{m+1}'$ are algebraically independent (see~\cite[(2.12)]{Mac}).
If $i\le m+1$, then $F^{(i)}(y)$ is a~polynomial at most of $y_1,\dots,y_{m+1}$.
Thus $F^{(i)}\big({-}\frac{y'}{2}\big)=0$ implies $F_\gamma=0$ for any $\gamma$ satisfying ${\rm wt}\,\gamma=i$.
Since $m$ is arbitrary we have $F_\gamma=0$ for any $\gamma$.
\end{proof}
